% ================================================================
% ==== FILE: content/m_spaces/m3.tex
% ================================================================

% ==============================
%  Ontology of Continua — Core
%  Meta-Space: M3
%  FULL MODULE — FINAL VERSION
% ==============================

\subsubsection{\texorpdfstring{$M_3$}{M_3} Übersicht}
\label{sec:m3-overview}

$M_3$ ist der Metaraum, der die Zulässigkeitsbedingungen dafür bereitstellt
dreidimensionale Kontinua $K_3$.
Es ist der erste Metaraum, in dem:

\begin{itemize}
    \item wird eine dritte unabhängige Achse $A_3$ zulässig,
    \item reaktive Struktur erscheint (chemische Bindung / Dissoziation),
    \item-Valenzzustände werden Teil des zulässigen Konfigurationsraums,
    \item-Diffusion und -Transport erfolgen in drei Raumdimensionen,
    \item lokale Interaktionsgraphen (Reaktionsnetzwerke) sind stabile Objekte,
    \item Aktivierungsschwellen und Reaktionspotentiale werden klar definiert,
    \item räumliche Heterogenität und Gradienten werden dynamisch unterstützt.
\end{itemize}

$M_3$ ist die mathematische und physikalische Umgebung, in der
Das chemische Kontinuum $K_3$ kann existieren, sich entwickeln und strukturelle Übergänge durchlaufen
in Richtung höherer Kontinua wie $K_4$.

% ---------------------------------------------------------------
\subsubsection*{1. Zulässiger Konfigurationsraum \texorpdfstring{$\Omega(M_3)$}{\Omega(M_3)}}

Der zulässige Satz von Konfigurationen in $M_3$ ist:

\[
\Omega(M_3)
=
\Big\{
(\phi,\, G_{\mathrm{chem}},\, \rho,\, g_{ij}) \ \Big|\
\phi \in C^{0}(X_3, V),\
G_{\mathrm{chem}} \in \mathcal{G}_{\mathrm{RAF}},\
\rho \in L^1(X_3),\
g_{ij}\in C^{0}
\Big\}.
\]

Wesentliche Zulassungsbedingungen:

\begin{enumerate}
    \item \textbf{Dreidimensionale Topologie.}
    $X_3$ muss eine 3-Mannigfaltigkeit sein, die lokale Koordinatendiagramme unterstützt.
    Differenzierbarkeit und Diffusion.

    \item \textbf{Lokale Reaktionsnetzwerke.}
    $G_{\mathrm{chem}}$ muss als endliches, lokal unterstütztes Element darstellbar sein
    Reaktionsdiagramm, das die RAF-Bedingungen für Schließung und katalytische Wirkung erfüllt.

    \item \textbf{Dichtefelder.}
    $\rho(x)$ definiert molekulare oder atomare Konzentrationen
    und muss mit endlicher Gesamtmasse integrierbar sein.

    \item \textbf{Metrikzulässigkeit.}
    $g_{ij}$ beschreibt räumliche Geometrie, schließt sie aber noch nicht ein
    relativistische Struktur; nur Riemannsche (positiv definite) Metriken
    sind auf dieser Ebene zulässig.

    \item \textbf{Ort.}
    Alle Interaktionen müssen durch lokale Potenziale ausdrückbar sein
    und lokale Reaktionskerne.
\end{enumerate}

Somit erweitert $\Omega(M_3)$ die geometrische und physikalische Zulässigkeit von $M_2$
in das chemische Regime.

% ---------------------------------------------------------------
\subsubsection*{2. Zulässige Achsen \texorpdfstring{$A(M_3)$}{A(M_3)}}

Eine neue Achse $A_3$ wird in $M_3$ zulässig und erfüllt:

\[
A_3 \not\subset \operatorname{span}(A_1, A_2),
\qquad
T_2 \ge\Theta_{\mathrm{dim}}.
\]

Zusätzlich zulässige interne Achsen:

\begin{itemize}
    \item $A_{\mathrm{val}}$: Valenzzustände (Bindungsmöglichkeiten),
    \item $A_{\mathrm{react}}$: zulässige Reaktionsmodi,
    \item $A_{\mathrm{conf}}$: Konfigurationszustände von Molekülen.
\end{itemize}

Diese internen Achsen kodieren die chemische Struktur innerhalb des Metaraums.

% ---------------------------------------------------------------
\subsubsection*{3. Potentiale \texorpdfstring{$P(M_3)$}{P(M_3)}}

Zulässige Potentiale in $M_3$ entsprechen chemischen und räumlichen Wechselwirkungen:

\[
P(M_3)
=
\left\{
U_{\mathrm{bond}} + U_{\mathrm{rep}} + U_{\mathrm{conf}}
+ U_{\mathrm{diff}} + U_{\mathrm{local}}
\right\}.
\]

Wo:

\begin{itemize}
    \item $U_{\mathrm{bond}}$ — Bindungspotentiale mit Minima bei stabilen Bindungslängen,
    \item $U_{\mathrm{rep}}$ – Abstoßungspotentiale mit kurzer Reichweite, die einen Zusammenbruch verhindern,
    \item $U_{\mathrm{conf}}$ — interne Konfigurationspotentiale für molekulare Zustände,
    \item $U_{\mathrm{diff}}$ — Potentiale, die mit Konzentrationsgradienten verbunden sind,
    \item $U_{\mathrm{local}}$ – jeder lokal definierte energetische Beitrag.
\end{itemize}

Einschränkungen:

\begin{enumerate}
    \item-Potentiale müssen nach unten begrenzt werden,
    \item sie müssen fast überall differenzierbar sein,
    \item Sie müssen eine lokale und keine globale Struktur oder eine Aktion aus der Ferne erzwingen.
\end{enumerate}

% ---------------------------------------------------------------
\subsubsection*{4. Schwellenwerte \texorpdfstring{$\Theta(M_3)$}{\Theta(M_3)}}

$M_3$ formalisiert mehrere chemische Schwellenwerte:

\begin{itemize}
    \item $\Theta_{\mathrm{bond}}$ – minimale Energie, die zum Aufbrechen einer Bindung erforderlich ist,
    \item $\Theta_{\mathrm{act}}$ — Aktivierungsschwelle für Reaktionen,
    \item $\Theta_{\mathrm{cluster}}$ — minimale stabile Clusterenergie,
    \item $\Theta_{\mathrm{grad}}$ – Schwelle für die Aufrechterhaltung räumlicher Gradienten,
    \item $\Theta_{\mathrm{RAF}}$ – Schwelle für die Schließung des RAF-Netzwerks.
\end{itemize}

Stabilität von $K_3$ erfordert:

\[
T_3 < \min\{\Theta_{\mathrm{bond}}, \Theta_{\mathrm{cluster}},
\Theta_{\mathrm{RAF}}\}.
\]

Wobei $T_3$ die strukturelle Spannung von Konfigurationen in $\Omega(K_3)$ ist.

% ---------------------------------------------------------------
\subsubsection*{5. Flüsse \texorpdfstring{$J(M_3)$}{J(M_3)}}

Zulässige Flüsse umfassen:

\[
J(M_3) =
\big(
J_{\mathrm{diff}},\
J_{\mathrm{reagieren}},\
J_{\mathrm{grad}},\
\nabla_i T_3
\Big).
\]

Interpretation:

\begin{itemize}
    \item $J_{\mathrm{diff}}$ — Diffusionsströmungen in 3D,
    \item $J_{\mathrm{react}}$ – lokale Reaktionsflüsse bei gegebenen RAF-Regeln,
    \item $J_{\mathrm{grad}}$ – Strömungen, die lokale Gradienten aufrechterhalten oder reduzieren,
    \item $\nabla_i T_3$ – Strömungen, die durch Gradienten struktureller Spannung angetrieben werden.
\end{itemize}

Alle Ströme müssen die lokale Massenerhaltung wahren.

% ---------------------------------------------------------------
\subsubsection*{6. Zyklen und Reaktionsnetzwerke}

$M_3$ ist der erste Metaraum, in dem chemische Kreisläufe zulässig sind:

\[
C_{\mathrm{chem}}
=
\big\{
\text{Bindungsbildung} \to
\text{Transformation} \to
\text{Bindungsbruch}
\to \dots
\Big\}.
\]

RAF-Netze werden zu zulässigen Zyklen, wenn:

\[
E_{\mathrm{in}} \ge \Theta_{\mathrm{act}}
\quad\text{und}\quad
\text{Abschlussbedingungen erfüllt}.
\]

Diese Zyklen bilden die minimalen Bausteine für die Entstehung biologischer Kontinua ($K_4$).

% ---------------------------------------------------------------
\subsubsection*{7. Grenze \texorpdfstring{$\partial\Omega(M_3)$}{\partial\Omega(M_3)} und Kollaps}

Eine Konfiguration nähert sich der Grenze $\partial\Omega(M_3)$, wenn:

\begin{itemize}
    \item Bindungspotentiale werden singulär,
    \item Reaktionsfluss divergiert,
    \item-Gradienten werden nicht integrierbar ($|\nabla\rho| \to \infty$),
    \item RAF-Abschluss geht verloren ($G_{\mathrm{chem}} \notin \mathcal{G}_{\mathrm{RAF}}$),
    \item Gesamtstrukturspannung überschreitet Schwellenwert:
    \[
    T_3 \ge \Theta_3.
    \]
\end{itemize}

Das Überqueren von $\partial\Omega(M_3)$ zerstört die chemische Kontinuität und kollabiert $K_3$.

% ---------------------------------------------------------------
\subsubsection*{8. Beziehung zu \texorpdfstring{$M_2$}{M_2} und \texorpdfstring{$M_4$}{M_4}}

Der Übergang $K_2 \to K_3$ wird zulässig, wenn:

\[
A_3 \in A(M_3),
\qquad
\Omega(K_3) \subseteq \Omega(M_3),
\qquad
T_2 \ge\Theta_{\mathrm{dim}}.
\]

Der Übergang zu $K_4$ erfordert:

\[
\Omega(K_4) \subseteq \Omega(M_4),
\qquad
\partial\Omega(M_3) \ \text{unterstützt begrenzte Permeabilität},
\qquad
\Theta_{\mathrm{mem}} > 0.
\]

Somit ist $M_3$ der Metaraum von:

\begin{itemize}
    \item chemische Struktur,
    \item Reaktionsnetzwerke,
    \item dreidimensionale Diffusion,
    \item Valenz und Bindung,
    \item Aktivierungsbarrieren,
    \item robuste Farbverläufe.
\end{itemize}

Es ist die strukturelle Grundlage, die chemische Kontinua $K_3$ und die Entstehung protozellulärer Strukturen in $K_4$ ermöglicht.
